% Layout-only edition, 22 September 2026.
% Controlling mathematical source: editorial/solver_dossiers_20260911/p121_reduced_excellent_curve_links_20260911/CONDUCTOR_IDEALS.tex
% Source SHA-256: 73eb8ce8754222fda5623965f92a6eeefb536a33d5949109c8387aaafd7e732f
% Only display grouping and line breaks differ; the source is retained.
% Book-local P12.2/P12.3/P12.6 conductor-controlled curve-link result.
% This is an ideal-function theorem on a finite normalized link, not a
% construction of the full singular local psh or Monge--Ampere theory.
\begin{proposition}[Conductor realization of branch orders]
\label{prop:p126-curve-conductor-principal-realization}
Let \((R,\mathfrak m)\) be a one-dimensional reduced Nagata noetherian
local ring, put \(D=R^\nu\), and enumerate
\[
             \operatorname{Br}(R,\mathfrak m)
                     =\{\mathfrak q_1,\ldots,\mathfrak q_s\}.
\]
Let
\[
 \mathfrak c=(R:D):=\{b\in D:bD\subseteq R\}
\]
be the conductor.  For every \(i\), define \(c_i\geq0\) by
\[
 \mathfrak cD_{\mathfrak q_i}
     =\mathfrak q_i^{c_i}D_{\mathfrak q_i},
 \qquad C_i:=\max\{c_i,1\}.
 \tag{\(\clubsuit\)}\label{eq:p126-curve-conductor-orders}
\]
Then, for every tuple of integers \(k_i\geq C_i\), there is an
\(R\)-nonzerodivisor \(f\in\mathfrak m\) such that
\[
                              w_{\mathfrak q_i}(f)=k_i
                       \qquad(1\leq i\leq s).
 \tag{\(\spadesuit\)}\label{eq:p126-prescribed-branch-orders}
\]
In particular, \((f)\) is an \(\mathfrak m\)-primary principal ideal.
Thus every sufficiently deep integral effective divisor on the finite
normalization fibre is cut out by one nonzerodivisor downstairs.
\end{proposition}

\begin{proof}
By \Cref{thm:p121-reduced-nagata-curve-link}, the finite \(R\)-algebra
\(D\) is a semilocal finite product of Dedekind domains and the displayed
\(\mathfrak q_i\) are all of its maximal ideals.  The conductor contains a
nonzerodivisor.  Indeed, choose finite \(R\)-module generators
\(b_1,\ldots,b_t\) of \(D\).  Since each \(b_j\) lies in the total ring of
fractions of \(R\), there is an \(R\)-nonzerodivisor \(a_j\) with
\(a_jb_j\in R\).  Their product \(a=\prod_ja_j\) is again a
nonzerodivisor and satisfies \(aD\subseteq R\), so \(a\in\mathfrak c\).

Consequently \(\mathfrak c\) is a nonzero ideal on every Dedekind factor.
Unique ideal factorization gives
\[
                         \mathfrak c
                    =\prod_{i=1}^s\mathfrak q_i^{c_i}
\]
with the exponents in \eqref{eq:p126-curve-conductor-orders}; see
\cite[Tags~034W and~034X]{StacksProject}.  For the prescribed tuple put
\[
                              J=\prod_{i=1}^s\mathfrak q_i^{k_i}.
\]
Then \(J\subseteq\mathfrak c\subseteq R\).  The ideal \(J\) is invertible,
hence each \(J/\mathfrak q_iJ\) is a nonzero one-dimensional vector space
over \(D/\mathfrak q_i\).  The Chinese remainder theorem for the pairwise
comaximal maximal ideals \(\mathfrak q_i\) lets us choose \(f\in J\) whose
class is nonzero in every \(J/\mathfrak q_iJ\).  Nakayama's lemma then shows
that \(f\) generates \(J_{\mathfrak q_i}\) over \(D_{\mathfrak q_i}\) for
every \(i\).  These are all the maximal ideals of \(D\), so \(J/fD=0\) and
\(J=fD\).  The inclusion \(J\subseteq R\) gives \(f\in R\), while
localization at \(\mathfrak q_i\) gives
\[
                         fD_{\mathfrak q_i}
                    =\mathfrak q_i^{k_i}D_{\mathfrak q_i},
\]
which is \eqref{eq:p126-prescribed-branch-orders}.

Multiplication by this basis element identifies \(D\) with \(J\), so \(f\)
is a nonzerodivisor of \(D\), and hence of \(R\subseteq D\).  Since every
\(k_i\geq1\), the element \(f\) lies in every \(\mathfrak q_i\), and thus
in \(\mathfrak m\).  A prime of the one-dimensional local ring \(R\)
which contains \(f\) cannot be minimal, because \(f\) is a
nonzerodivisor.  It must therefore be \(\mathfrak m\).  Hence
\(\sqrt{(f)}=\mathfrak m\); a noetherian ideal with maximal radical is
primary.
\end{proof}

\begin{theorem}[Principal-ideal profiles and envelopes on a curve link]
\label{thm:p122-curve-principal-ideal-envelope}
In the setting of
\Cref{prop:p126-curve-conductor-principal-realization}, let
\(\mathfrak a\subset R\) be \(\mathfrak m\)-primary and write
\[
 d_i=w_{\mathfrak q_i}(\mathfrak a),
 \qquad v_i=\frac{w_{\mathfrak q_i}}{d_i}
       \in\mathcal V_{\mathfrak a}(R).
\]
For an \(R\)-nonzerodivisor \(f\) such that \((f)\) is
\(\mathfrak m\)-primary and an integer \(N\geq1\), put
\[
 u_{f,N}(v):=\frac1N\log|(f)|(v)=-\frac{v(f)}N,
 \qquad
 \mathscr P_{\mathfrak a}^{\mathrm{pr}}
 :=\{u_{f,N}:f,N\text{ as above}\}.
 \tag{\(\diamondsuit\)}\label{eq:p122-principal-ideal-profile-family}
\]
Then the following assertions hold.
\begin{enumerate}
\item The uniform closure of this explicitly defined family is the complete
      nonpositive orthant.  More precisely,
      \[
       \begin{gathered}
       \mathscr P_{\mathfrak a}^{\mathrm{pr}}
       =\{u:u(v_i)\in\mathbf Q_{<0}\text{ for every }i\},\\
       \qquad
       \overline{\mathscr P_{\mathfrak a}^{\mathrm{pr}}}^{\,\|\cdot\|_\infty}
       =\{u:\mathcal V_{\mathfrak a}(R)\to\mathbf R:u\leq0\}.
       \end{gathered}
       \tag{\(\heartsuit\)}\label{eq:p122-principal-ideal-orthant}
      \]
\item Every nonpositive real-valued function is the increasing uniform limit
      of a sequence in \(\mathscr P_{\mathfrak a}^{\mathrm{pr}}\).  Every
      strictly negative real-valued function is also the decreasing uniform
      limit of such a sequence.  If \(u(v_i)=0\) for some branch, no
      decreasing sequence in \(\mathscr P_{\mathfrak a}^{\mathrm{pr}}\)
      can converge pointwise to \(u\).
\item For every finite real-valued obstacle \(\phi\) on the link, define the
      principal-ideal envelope by
      \[
       P_{\mathfrak a}^{\mathrm{pr}}\phi(v)
       :=\sup\{u(v):u\in\mathscr P_{\mathfrak a}^{\mathrm{pr}},\ u\leq\phi\}.
       \tag{\(\star\)}\label{eq:p123-principal-ideal-envelope}
      \]
      Then
      \[
                         P_{\mathfrak a}^{\mathrm{pr}}\phi
                              =\min\{\phi,0\}.
       \tag{\(\dagger\)}\label{eq:p123-principal-ideal-envelope-formula}
      \]
      The supremum in \eqref{eq:p123-principal-ideal-envelope} admits an
      increasing uniformly convergent sequence of principal-ideal
      competitors.  The same formula results if the competitor family is
      replaced by its uniform closure in
      \eqref{eq:p122-principal-ideal-orthant}.  The original supremum need
      not be attained by a member of
      \(\mathscr P_{\mathfrak a}^{\mathrm{pr}}\).
\end{enumerate}
These are statements about the displayed ideal-function family on a finite
link.  They do not identify that family with an independently constructed
singular nef or psh cone and do not construct a Monge--Ampere product or an
energy class.
\end{theorem}

\begin{proof}
Write \(u(v_i)=-\alpha_i\), where \(\alpha_i\geq0\).  With \(C_i\) as in
\eqref{eq:p126-curve-conductor-orders}, set
\[
 k_{n,i}=\left\lceil2^nd_i\alpha_i\right\rceil+C_i.
 \]
The conductor realization proposition supplies \(f_n\) with these branch
orders.  For \(N=2^n\),
\[
 u_{f_n,2^n}(v_i)=-\frac{k_{n,i}}{2^nd_i}\leq-\alpha_i=u(v_i),
 \qquad
 0\leq u(v_i)-u_{f_n,2^n}(v_i)
       <\frac{C_i+1}{2^nd_i}.
 \tag{\(\ddagger\)}\label{eq:p122-increasing-principal-approximation}
\]
Since \(\lceil2x\rceil\leq2\lceil x\rceil\), one has
\(k_{n+1,i}\leq2k_{n,i}\).  Thus the functions in
\eqref{eq:p122-increasing-principal-approximation} increase pointwise, and
the finite branch set makes the convergence uniform.

If every \(\alpha_i>0\), then for all sufficiently large \(n\) the integers
\[
                    \ell_{n,i}=\left\lfloor2^nd_i\alpha_i\right\rfloor
\]
satisfy \(\ell_{n,i}\geq C_i\).  Choose \(g_n\) with those branch orders.
Then
\[
 u_{g_n,2^n}(v_i)=-\frac{\ell_{n,i}}{2^nd_i}\geq u(v_i),
 \qquad
 0\leq u_{g_n,2^n}(v_i)-u(v_i)<\frac1{2^nd_i}.
\]
The inequality \(\lfloor2x\rfloor\geq2\lfloor x\rfloor\) shows that this
sequence decreases.  Conversely, every member of
\(\mathscr P_{\mathfrak a}^{\mathrm{pr}}\) is strictly negative at every
\(v_i\).  A decreasing sequence starting with such a function therefore
cannot have limit zero at any branch.

If all \(\alpha_i\) are positive rationals, choose one sufficiently
divisible \(N\) such that every \(k_i=Nd_i\alpha_i\) is an integer at least
\(C_i\).  A single application of the conductor proposition then gives
\(u=u_{f,N}\).  Conversely, every member of the displayed family has
strictly negative rational coordinates.  Together with the increasing
approximation, this proves
\eqref{eq:p122-principal-ideal-orthant} and the first two assertions.

Finally put \(w=\min\{\phi,0\}\).  The increasing construction applied to
\(w\) gives principal-ideal profiles \(u_n\leq w\leq\phi\) converging
uniformly to \(w\).  Hence the envelope is at least \(w\).  Every admissible
profile is both nonpositive and bounded above by \(\phi\), so it is bounded
above by \(w\); this proves
\eqref{eq:p123-principal-ideal-envelope-formula}.  The closure version is
identical by \eqref{eq:p122-principal-ideal-orthant}.
\end{proof}
